\section{Results}
\label{sec:results}

\paragraph{Smaller models are more biased.}
\textit{Evidence.} Across the cross-model sweep ($n{=}770$ per cell), abstention on
ambiguous items rises with scale within every family, three-option abstention goes
from $0.77$ to $0.89$ for Qwen ($0.6$\,B$\to$$4$\,B) and from $0.30$ to $0.89$ for
Llama ($1$\,B$\to$$3$\,B), and disambiguated accuracy rises with scale as well. The
smallest model, \texttt{Llama-3.2-1B}, abstains least ($0.30$, below the $1/3$ chance
line) and most often names the stereotyped target on disambiguated items ($0.68$ vs.\
$0.40$ for \texttt{Qwen3-0.6B}). In the bias-vs-accuracy view (\Cref{fig:results-bias}),
small models are both less accurate and more biased, while larger models cluster near
zero bias at high accuracy. All rates carry Wilson $95\%$ confidence intervals.\\
\textit{Limitation.} The sweep mixes families and sizes, so we do not cleanly resolve
size from family on every axis.\\
\textit{Notes.} One run (\texttt{gemma-2-27b-it}) was degenerate and is shown honestly
rather than dropped.

\plotfig{figures/baseline_bias_alignment_old_new.png}{baseline_bias_alignment_old_new.png}{0.98\linewidth}{\textbf{Bias alignment vs.\ accuracy} across both the old baseline models and the new sweep ($29$ models / $58$ segments, count-based $F(\text{Target})/F(\text{Other})$). Each model is a segment at its accuracy spanning its neutral-vs-negative wording bias ($x{<}0$ leans toward the Other group, $x{>}0$ toward the stereotyped Target). Left: ambiguous items (accuracy $=$ abstention); right: disambiguated. Small models are less accurate and more biased; large models cluster near zero bias at high accuracy.}{fig:results-bias}

\paragraph{Scale lowers uncertainty.}
\textit{Evidence.} We fork the chain of thought at every token and track the outcome
distribution $O_t$ (\Cref{fig:results-forking}). The larger \texttt{Qwen3-14B} commits
its answer \emph{decisively}: the change-point test finds a single significant
commitment token (at $t{=}270$, Bayes factor ${\approx}10^{9}$), after which the answer
snaps from abstention to a named group. The smaller \texttt{Qwen3-0.6B} instead
accretes its answer gradually, with no single deciding token (Bayes factor $0.023$).
The larger model resolves to a sharp, low-uncertainty commitment.\\
\textit{Limitation.} A single-item case study per model.\\
\textit{Notes.} The forking (red) token is the position where re-sampling a different
token most diverts the final outcome.

\begin{figure}[htbp]\centering
  \begin{subfigure}{0.49\linewidth}\centering
    \plotsub{../out/sesgo/forking/Qwen3-0.6B/plots/forking_commit_dynamics.png}{../out/sesgo/forking/Qwen3-0.6B/plots/forking_commit_dynamics.png}{\linewidth}
    \caption{\texttt{Qwen3-0.6B}: the answer wanders, with no single commitment token.}
  \end{subfigure}\hfill
  \begin{subfigure}{0.49\linewidth}\centering
    \plotsub{figures/forking_ot_14b.png}{forking_ot_14b.png}{\linewidth}
    \caption{\texttt{Qwen3-14B}: a single decisive commitment token ($t{=}270$, red).}
  \end{subfigure}
  \caption{\textbf{Scale lowers uncertainty.} Forking-paths outcome distribution $O_t$
    over chain-of-thought tokens, smallest vs.\ larger model (target / other / unknown /
    unparseable). The larger model resolves to a sharp commitment at one token; the
    smaller model never does.}
  \label{fig:results-forking}
\end{figure}

\paragraph{Smaller models do not answer consistently.}
\textit{Evidence.} We render each item into three orthogonal variants --- a
target/other position flip and a label-style swap, gold answer fixed --- and ask how
often the committed answer never changes, now across the \emph{full model sweep} rather
than a single smallest-vs-largest pair. Format-invariance rises steeply with scale
within every family: \texttt{Llama} from $0.23$ at $1$B to $0.60$ at $8$B and
\texttt{Qwen3.5} from $0.54$ at $0.8$B to $0.97$ at $27$B (\Cref{fig:results-stability}).
A purely cosmetic relabelling flips the small models far more often than the large.\\
\textit{Limitation.} Coverage is partial on the largest slices (small $n$, annotated).\\
\textit{Notes.} This implies the other studies should filter for unstable answers.

\plotfig{figures/new_stability_format_invariance.png}{new_stability_format_invariance.png}{0.95\linewidth}{\textbf{Smaller models do not answer consistently.} Fraction of items whose committed role is identical across all three orthogonal format variants, across the full model sweep (Wilson $95\%$ CIs, $n$ items annotated; filled triangles thinking, open circles non-thinking; faint markers are order-only and label-only agreement). Format-invariance rises sharply with scale within every family --- scale, not family or formatting, governs answer stability.}{fig:results-stability}
